\chapter{Objectives}
\label{ch:objectives}

\section{General Objective}
The general objective of this project is to build a self-guided,
geofence-verified AR exploration platform that makes heritage and campus
visits engaging, measurable, and scalable --- replacing the human-guided
tour with software that guides visitors in real time, verifies their
physical presence, and rewards exploration.

\section{Specific Objectives}
To realise the general objective, \projectname{} pursues the following
specific objectives:

\begin{itemize}
  \item \textbf{Verify physical presence via dynamic geofencing.} Detect and
        confirm a visitor's arrival at each waypoint using a geofence radius
        of 10~metres stored per place, so that visits, rewards, and
        analytics are grounded in verified presence rather than
        self-reporting.
  \item \textbf{Guide visitors in real time.} Provide continuous wayfinding
        to the next waypoint through a live compass heads-up display with a
        GPS distance readout, together with an AR camera overlay that
        anchors guidance in the visitor's view of the real world.
  \item \textbf{Unlock location-specific media on confirmed arrival.}
        Release each checkpoint's images and videos --- historical context,
        testimonials, and narrative clips --- only after the geofence has
        confirmed that the visitor is physically at the location.
  \item \textbf{Drive engagement through gamification.} Sustain exploration
        with AR Points, milestone completions, optional side quests,
        post-route knowledge quizzes, digital badges, a community
        leaderboard, and a store where points are redeemed for tangible
        perks.
  \item \textbf{Give administrators no-code control.} Allow site
        administrators to create and manage routes, waypoints, media,
        quizzes, and users entirely through a dashboard, with engagement
        analytics generated organically from per-place visitor counts.
  \item \textbf{Remain hardware-free and site-agnostic.} Run entirely in the
        smartphone browser with no physical installation at the site, so
        that the same platform can be configured for any campus, museum, or
        heritage property.
\end{itemize}
